% --- PACOTES ---
%> Pacotes essenciais
\usepackage[
    a0paper,
    top=2cm,
    bottom=2cm,
    left=2cm,
    right=2cm
]{geometry}
\usepackage[
    backend=biber,
    style=numeric-comp,
    sorting=nyt,
    natbib=true,
    maxcitenames=1,
    mincitenames=1,
    maxbibnames=1,
    minbibnames=1,
    url=false, doi=false, isbn=false, eprint=false
]{biblatex}
\renewcommand*{\bibfont}{\setBibFont}
\setlength{\bibitemsep}{0.5em}
\addbibresource{_bibliography/references.bib}
\usepackage{graphicx}
\usepackage{indentfirst}
\usepackage[T1]{fontenc}
\usepackage[brazilian]{babel}

%> Pacotes de estilo
\usepackage{anyfontsize}
\newcommand*{\ClipSep}{0.4cm}%
\usepackage{fontawesome5}
\usepackage{multicol}
\setlength{\columnsep}{3em}
\usepackage{caption}
\usepackage{microtype}
\usepackage[autostyle]{csquotes}
\usepackage{booktabs}
\usepackage{enumitem}
\usepackage{array}
\renewcommand{\arraystretch}{1.6}
\usepackage[svgnames]{xcolor}
\usepackage{qrcode}

%> Pacotes matemáticos
\usepackage{amsmath, amssymb, amsfonts}
\usepackage{esint}
\usepackage{upgreek}
\usepackage{tikz}
\usetikzlibrary{angles, quotes, arrows.meta, babel}
\usepackage{newpxtext, newpxmath}
\usepackage[most]{tcolorbox}

%> Set fontes
\DeclareCaptionFont{basefont}{\setBaseFont}
\captionsetup{font=basefont}

%> Gerador de blá blá blá
\newcommand{\blablabla}[1][1]{%
  \foreach \n in {1,...,#1}{%
    Neste parágrafo vamos simular um texto em português para testar a tipografia, o espaçamento, a hifenização e a formatação geral do documento. Como estamos desenvolvendo um modelo robusto para a Universidade Federal de Viçosa (UFV), é importante que as margens sejam testadas com palavras reais da nossa língua, contendo acentuações, cedilhas e construções frasais típicas do nosso idioma. Dessa forma, garantimos que o código do template funcione perfeitamente, sem depender de pacotes externos ou línguas mortas. \par\vspace{0.2cm}%
  }%
}